\pg{334} 56. \textbf{O ente, que é acto puro, é necessario ; e o ente, que é necessario, é acto puro.}

\textbf{a)} \emph{O ente, que é acto puro, é necessario.} --- O ente, que é \emph{acto puro,} exclue de si toda a potencia; aliás não seria o que é. Ora o ente, que exclue de si toda a potencia, deve possuir a existencia identificada com a essencia, e por isso é \emph{necessario.} Logo o ente, que é \emph{acto puro,} é \emph{necessario}\footnote{O \emph{acto puro} exclue de si toda a potencia. Ora o ente, que exclue de si toda a potencia, não pode ser \emph{contingente;} porque o ente \emph{contingente} não exclue, mas encerra a potencia, emquanto está em potencia para o acto da existencia. Portanto o ente, que é \emph{acto puro,} é \emph{necessario,} possue a existencia pela exigencia da sua essencia. Nem podia ser diversamente. O ente, que é \emph{acto puro,} possue a existencia identificada com a essencia, de modo que podemos dizer que a existencia é um dos elementos da essencia. Sendo assim, esse ente é \emph{necessariamente existente,} como o homem, que pela sua essencia é dotado de razão, é \emph{necessariamente racional.} As essencias, como vimos, são \emph{necessarias} quanto aos elementos, que as constituem.}.

\textbf{b)} \emph{O ente, que é necessario, é acto puro.} --- O ente, que é \emph{necessario,} existe pela exigencia da sua essencia, e por isso possue a existencia identificada com a essencia. Ora o ente, que possue a existencia identificada com a essencia, é \emph{acto puro.} Logo o ente, que é \emph{necessario,} é \emph{acto puro}\footnote{O ente \emph{necessario} é \emph{acto puro}: O \emph{acto} é \emph{puro,} quando não é recebido num sujeito distincto, nem está em potencia para um acto ulterior. Ora o ente \emph{necessario} --- não é recebido num sujeito distincto, porque o seu \emph{ser} é uma e a mesma coisa com a \emph{essencia,} --- nem está em potencia para um acto ulterior, porque o ente, cujo ser se identifica com a essencia, é o proprio ser, é o ser subsistente, e por isso possue a plenitude da perfeição.\par D' onde se vê que o ente \emph{necessario} é \emph{unico.} Por quanto o \emph{acto,} como dissemos, só se multiplica emquanto é recebido em muitos sujeitos. Ora o ser do ente \emph{necessario,} identificando-se com a essencia, não é recebido em nenhum sujeito, mas é por si subsistente, e por isso não se multiplica, mas é \emph{unico.} Logo o ente \emph{necessario} é \emph{unico.}}.

57. \textbf{O ente, que consta de acto e de potencia, é contingente; e o ente, que é contingente, consta de acto e de potencia.}

\textbf{a)} \emph{O ente, que consta de acto e de potencia, é contingente.} --- O ente, que consta de \emph{acto} e de \emph{potencia,} possue a existencia \emph{realmente} distincta da essencia. Ora o ente, que possue a existencia \emph{realmente} distincta da essencia, é \emph{contingente.} Logo o ente, que consta de \emph{acto} e de \emph{potencia,} é \emph{contingente}\footnote{A conclusão é certissima. Na verdade, todo o ente composto pode ser dissolvido; e, podendo ser dissolvido, é naturalmente \emph{contingente.} --- Além d' isso, o ente composto é \emph{posterior} ás partes componentes, e por isso não pode ser \emph{necessario,} mas deve ser \emph{contingente;} porque o ente \emph{necessario,}\pgn{335}sendo o proprio ser, não está ordenado para outro ser mais perfeito, e é simplesmente o \emph{primeiro.} --- Mais: o composto depende, no seu ser, de uma causa, que une as parte componentes e que deve ser não só anterior ao proprio composto, mas distincta delle. Dependendo, no seu ser, de uma causa distincta, o composto não existe pela sua natureza, e por isso é \emph{contingente.}}.